% GENERATED FILE. Edit canonical lessons, then run npm run generate:books.
% canonical-source-hash: fnv1a64:064fe22f4dce5b43
% canonical-lessons: LA-C388-municeps, LA-C388-decurio, LA-C388-colonia, LA-C388-sodalis, LA-C388-collegium, LA-R388-review-381-388

\chapter{Towns and Associations}
\label{ch:la-towns-and-associations}
\hlchaptermodality{388}

\begin{chapteropening}
\textbf{By the end of this chapter:} \emph{I can talk about a townsman and a town councillor, a colony, a comrade and an association.}

Complete the last lesson of chapter 388: I can talk about a townsman and a town councillor, a colony, a comrade and an association.
\end{chapteropening}

\section[mūniceps, mūnicipis]{mūniceps, mūnicipis — a citizen of a free town; a fellow townsman}

\label{lesson:LA-C388-municeps}



\begin{quote}
Before the new one: say the Latin for stormy, then the Latin for to stir up.
\end{quote}

\subsection*{You'll want to know: mūniceps, mūnicipis}
\textbf{mūniceps, mūnicipis} — \textquotedblleft{}a citizen of a free town; a fellow townsman\textquotedblright{}.

\begin{quote}
Ille mūniceps meus est.
\end{quote}

\begin{quote}
\emph{He comes from my home town.}
\end{quote}

\emph{mūniceps, mūnicipis} m. and f.: one who took on the duties of a self-governing Italian town.

\begin{cousinweb}
\emph{mūniceps} is built from \emph{mūnus} \textquotedblleft{}a duty, a service\textquotedblright{} and \emph{capiō} \textquotedblleft{}to take\textquotedblright{} (EtymDB; Lewis \& Short): one who takes on the town's duties.
\end{cousinweb}

\begin{grammarlens}[title={What you've built}]
One more word for telling, explaining and describing at B1.
\end{grammarlens}

\subsection*{Your turn}
\begin{itemize}
\raggedright
  \item \emph{Say it:} \emph{mūniceps, mūnicipis}
  \item \emph{Say it:} \emph{mūniceps, mūnicipis}, once more
  \item \emph{Say it:} a sentence of your own with \emph{mūniceps, mūnicipis}
  \item \emph{Recall:} say the Latin for stormy, then the Latin for to stir up, then say \emph{mūniceps, mūnicipis} again
\end{itemize}

\begin{tcolorbox}[breakable,colback=teal!4,colframe=teal!35!black,title={Before you move on}]
What does mūniceps, mūnicipis mean? (\textquotedblleft{}A citizen of a free town; a fellow townsman\textquotedblright{}.) Say it once more.
\end{tcolorbox}
\section[decuriō, decuriōnis]{decuriō, decuriōnis — a town councillor; a cavalry officer}

\label{lesson:LA-C388-decurio}



\begin{quote}
Before the new one: say the Latin for to stir up, then the Latin for a citizen of a free town.
\end{quote}

\subsection*{You'll want to know: decuriō, decuriōnis}
\textbf{decuriō, decuriōnis} — \textquotedblleft{}a town councillor; a cavalry officer\textquotedblright{}.

\begin{quote}
Decuriōnēs dē aquaeductū dēlīberāvērunt.
\end{quote}

\begin{quote}
\emph{The town councillors deliberated about the aqueduct.}
\end{quote}

\emph{decuriō, decuriōnis} m. In a town the \emph{decuriōnēs} formed the council; in the army a \emph{decuriō} led ten horsemen. Built on \emph{decem}, ten.

\begin{cousinweb}
\emph{decuriō} is built on \emph{decuria} \textquotedblleft{}a group of ten\textquotedblright{} (EtymDB; Lewis \& Short), which comes from \emph{decem} \textquotedblleft{}ten\textquotedblright{}.
\end{cousinweb}

\begin{grammarlens}[title={What you've built}]
One more word for telling, explaining and describing at B1.
\end{grammarlens}

\subsection*{Your turn}
\begin{itemize}
\raggedright
  \item \emph{Say it:} \emph{decuriō, decuriōnis}
  \item \emph{Say it:} \emph{decuriō, decuriōnis}, once more
  \item \emph{Say it:} a sentence of your own with \emph{decuriō, decuriōnis}
  \item \emph{Recall:} say the Latin for to stir up, then the Latin for a citizen of a free town, then say \emph{decuriō, decuriōnis} again
\end{itemize}

\begin{tcolorbox}[breakable,colback=teal!4,colframe=teal!35!black,title={Before you move on}]
What does decuriō, decuriōnis mean? (\textquotedblleft{}A town councillor; a cavalry officer\textquotedblright{}.) Say it once more.
\end{tcolorbox}
\section[colōnia, colōniae]{colōnia, colōniae — a colony, a settlement}

\label{lesson:LA-C388-colonia}



\begin{quote}
Before the new one: say the Latin for a citizen of a free town, then the Latin for a town councillor.
\end{quote}

\subsection*{You'll want to know: colōnia, colōniae}
\textbf{colōnia, colōniae} — \textquotedblleft{}a colony, a settlement\textquotedblright{}.

\begin{quote}
Mīlitēs in colōniam missī sunt.
\end{quote}

\begin{quote}
\emph{Soldiers were sent to the settlement.}
\end{quote}

\emph{colōnia, colōniae} f., built on \emph{colōnus}, a settler, a farmer. Many cities of Europe began as Roman colonies. English \emph{colony}.

\begin{cousinweb}
\emph{colōnia} is built on \emph{colōnus} \textquotedblleft{}a settler, a farmer\textquotedblright{} (EtymDB), which comes from \emph{colō} \textquotedblleft{}to till, to dwell\textquotedblright{} (Lewis \& Short). English \emph{colony} comes from it.
\end{cousinweb}

\begin{grammarlens}[title={What you've built}]
One more word for telling, explaining and describing at B1.
\end{grammarlens}

\subsection*{Your turn}
\begin{itemize}
\raggedright
  \item \emph{Say it:} \emph{colōnia, colōniae}
  \item \emph{Say it:} \emph{colōnia, colōniae}, once more
  \item \emph{Say it:} a sentence of your own with \emph{colōnia, colōniae}
  \item \emph{Recall:} say the Latin for a citizen of a free town, then the Latin for a town councillor, then say \emph{colōnia, colōniae} again
\end{itemize}

\begin{tcolorbox}[breakable,colback=teal!4,colframe=teal!35!black,title={Before you move on}]
What does colōnia, colōniae mean? (\textquotedblleft{}A colony, a settlement\textquotedblright{}.) Say it once more.
\end{tcolorbox}
\section[sodālis, sodālis]{sodālis, sodālis — a companion, a comrade; a club member}

\label{lesson:LA-C388-sodalis}



\begin{quote}
Before the new one: say the Latin for a town councillor, then the Latin for a colony.
\end{quote}

\subsection*{You'll want to know: sodālis, sodālis}
\textbf{sodālis, sodālis} — \textquotedblleft{}a companion, a comrade; a club member\textquotedblright{}.

\begin{quote}
Cum sodālibus cēnāvit.
\end{quote}

\begin{quote}
\emph{He dined with his comrades.}
\end{quote}

\emph{sodālis, sodālis} m. and f. Members of a club, a guild or a priestly brotherhood were \emph{sodālēs}.

\begin{cousinweb}
The origin of \emph{sodālis} is not settled: Lewis \& Short compare it with a root meaning \textquotedblleft{}one's own way, custom\textquotedblright{}, but give no sure account.
\end{cousinweb}

\begin{grammarlens}[title={What you've built}]
One more word for telling, explaining and describing at B1.
\end{grammarlens}

\subsection*{Your turn}
\begin{itemize}
\raggedright
  \item \emph{Say it:} \emph{sodālis, sodālis}
  \item \emph{Say it:} \emph{sodālis, sodālis}, once more
  \item \emph{Say it:} a sentence of your own with \emph{sodālis, sodālis}
  \item \emph{Recall:} say the Latin for a town councillor, then the Latin for a colony, then say \emph{sodālis, sodālis} again
\end{itemize}

\begin{tcolorbox}[breakable,colback=teal!4,colframe=teal!35!black,title={Before you move on}]
What does sodālis, sodālis mean? (\textquotedblleft{}A companion, a comrade; a club member\textquotedblright{}.) Say it once more.
\end{tcolorbox}
\section[collēgium, collēgiī]{collēgium, collēgiī — an association, a guild; a board of officials}

\label{lesson:LA-C388-collegium}



\begin{quote}
Before the new one: say the Latin for a colony, then the Latin for a companion.
\end{quote}

\subsection*{You'll want to know: collēgium, collēgiī}
\textbf{collēgium, collēgiī} — \textquotedblleft{}an association, a guild; a board of officials\textquotedblright{}.

\begin{quote}
Fabrī in collēgium convēnērunt.
\end{quote}

\begin{quote}
\emph{The craftsmen met in their guild.}
\end{quote}

\emph{collēgium, collēgiī} n., built on \emph{collēga}, a colleague. Guilds of craftsmen and boards of priests were both \emph{collēgia}. English \emph{college}.

\begin{cousinweb}
\emph{collēgium} is built on \emph{collēga} \textquotedblleft{}a colleague, one chosen with another\textquotedblright{} (EtymDB; Lewis \& Short). English \emph{college} comes from it.
\end{cousinweb}

\begin{grammarlens}[title={What you've built}]
One more word for telling, explaining and describing at B1.
\end{grammarlens}

\subsection*{Your turn}
\begin{itemize}
\raggedright
  \item \emph{Say it:} \emph{collēgium, collēgiī}
  \item \emph{Say it:} \emph{collēgium, collēgiī}, once more
  \item \emph{Say it:} a sentence of your own with \emph{collēgium, collēgiī}
  \item \emph{Recall:} say the Latin for a colony, then the Latin for a companion, then say \emph{collēgium, collēgiī} again
\end{itemize}

\begin{tcolorbox}[breakable,colback=teal!4,colframe=teal!35!black,title={Before you move on}]
What does collēgium, collēgiī mean? (\textquotedblleft{}An association, a guild; a board of officials\textquotedblright{}.) Say it once more.
\end{tcolorbox}
\section[(dialogue)]{Words from chapters 381–388 — a review}

\label{lesson:LA-R388-review-381-388}



\begin{quote}
Say the Latin for an advocate and a client.
\end{quote}

\subsection*{Your turn}
Say these aloud:

\begin{itemize}
\raggedright
  \item an advocate, a client, a shipowner, a copyist, a land surveyor
  \item a senator, patrician, plebeian, well known, of the people
  \item a public office, a praetor, an aedile, a quaestor, a censor
  \item a vote, the voting assembly, a candidate, canvassing, to pick out
  \item to rule, to steer, to manage, command, the first
  \item a partnership, human nature, a custom, an institution, harmony
  \item disagreement, a rising, seditious, stormy, to stir up
  \item a citizen of a free town, a town councillor, a colony, a companion, an association
\end{itemize}

\begin{tcolorbox}[breakable,colback=teal!4,colframe=teal!35!black,title={Before you move on}]
An advocate? (\textbf{causidicus, causidicī}.) A client? (\textbf{cliēns, clientis}.) A shipowner? (\textbf{nāviculārius, nāviculāriī}.) A copyist? (\textbf{librārius, librāriī}.) A land surveyor? (\textbf{agrīmēnsor, agrīmēnsōris}.) A senator? (\textbf{senātor, senātōris}.) Patrician? (\textbf{patricius, patricia, patricium}.) Plebeian? (\textbf{plēbēius, plēbēia, plēbēium}.) Well known? (\textbf{nōbilis, nōbile}.) Of the people? (\textbf{populāris, populāre}.) A public office? (\textbf{magistrātus, magistrātūs}.) A praetor? (\textbf{praetor, praetōris}.) An aedile? (\textbf{aedīlis, aedīlis}.) A quaestor? (\textbf{quaestor, quaestōris}.) A censor? (\textbf{cēnsor, cēnsōris}.) A vote? (\textbf{suffrāgium, suffrāgiī}.) The voting assembly? (\textbf{comitia, comitiōrum}.) A candidate? (\textbf{candidātus, candidātī}.) Canvassing? (\textbf{ambitiō, ambitiōnis}.) To pick out? (\textbf{ēligō, ēligere}.) To rule? (\textbf{regō, regere}.) To steer? (\textbf{gubernō, gubernāre}.) To manage? (\textbf{administrō, administrāre}.) Command? (\textbf{imperium, imperiī}.) The first? (\textbf{prīnceps, prīncipis}.) A partnership? (\textbf{societās, societātis}.) Human nature? (\textbf{hūmānitās, hūmānitātis}.) A custom? (\textbf{mōs, mōris}.) An institution? (\textbf{īnstitūtum, īnstitūtī}.) Harmony? (\textbf{concordia, concordiae}.) Disagreement? (\textbf{discordia, discordiae}.) A rising? (\textbf{sēditiō, sēditiōnis}.) Seditious? (\textbf{sēditiōsus, sēditiōsa, sēditiōsum}.) Stormy? (\textbf{turbulentus, turbulenta, turbulentum}.) To stir up? (\textbf{concitō, concitāre}.) A citizen of a free town? (\textbf{mūniceps, mūnicipis}.) A town councillor? (\textbf{decuriō, decuriōnis}.) A colony? (\textbf{colōnia, colōniae}.) A companion? (\textbf{sodālis, sodālis}.) An association? (\textbf{collēgium, collēgiī}.)
\end{tcolorbox}
